
\documentclass{article}
\usepackage{colortbl}
\usepackage{makecell}
\usepackage{multirow}
\usepackage{supertabular}

\begin{document}

\newcounter{utterance}

\centering \large Interaction Transcript for game `wordle\_withclue', experiment `medium\_frequency\_words\_clue\_no\_critic', episode 20 with reflexion{-}Qwen3.5{-}9B.
\vspace{24pt}

{ \footnotesize  \setcounter{utterance}{1}
\setlength{\tabcolsep}{0pt}
\begin{supertabular}{c@{$\;$}|p{.15\linewidth}@{}p{.15\linewidth}p{.15\linewidth}p{.15\linewidth}p{.15\linewidth}p{.15\linewidth}}
    \# & \multicolumn{2}{c}{Player} && \multicolumn{2}{c}{Game Master} \\
    \hline

    \theutterance \stepcounter{utterance}  
    & & & \multicolumn{4}{p{0.6\linewidth}}{
        \cellcolor[rgb]{0.9,0.9,0.9}{
            \makecell[{{p{\linewidth}}}]{
                \texttt{\tiny{[P1$\langle$GM]}}
                \texttt{You have attempted this same task before and did not succeed. These are the plans you wrote after those attempts.} \\
\texttt{Trial \#1: New plan: I failed because my explanation became a verbose, rambling monologue that violated the "concise" constraint and likely exceeded the token limit or format expectations, causing the Game Master to reject the message. I also failed to actually provide a valid "guess" word in the final turn, only analyzing the logic. Next time, I will strictly adhere to the format: provide a single, concise sentence in the "explanation:" tag that explicitly references both the clue and the previous feedback, followed immediately by a valid five{-}letter word in the "guess:" tag. I will stop reasoning aloud and ensure the output contains exactly two tags with no extra text.} \\
\\ 
\texttt{Trial \#2: New plan: I failed because I got stuck in an infinite loop, repeatedly guessing "loons" after realizing it was incorrect, instead of exploring other valid five{-}letter winter wear words containing 'O' in the second position and 'L' elsewhere (such as "slosh" or "molds" if they fit the theme, or simply pivoting to a different strategy like testing common vowels/consonants). Next time, upon receiving negative feedback, I will immediately generate a new, distinct candidate word that satisfies the current constraints rather than re{-}submitting the same failed guess. I will also ensure my explanation explicitly mentions the specific letters to keep and discard based on the latest feedback before proposing the new word.} \\
\\ 
\texttt{Trial \#3: New plan: I failed because my explanation became verbose and rambling, violating the "concise" constraint, and I incorrectly deduced the target word structure (assuming a single 'o' when the feedback showed two 'o's with one green and one red, implying the word has two 'o's). Next time, I will strictly adhere to the format by providing a single, concise sentence in the "explanation:" tag that explicitly references both the clue ("winter wear") and the specific feedback from the previous turn (keeping 'o' and 'l', discarding 'w', 's', 'f', 'e', 'c'), followed immediately by a valid five{-}letter word in the "guess:" tag that satisfies these constraints, such as "coats" (if 'c' were valid) or a better fit like "wools" adjusted for the 'l' position or a new candidate like "snood" or "scarf" if they fit the letter constraints, ensuring the output contains exactly two tags with no extra text.} \\
\\ 
\\ 
\texttt{You are a language wizard who likes to guess words by using the given rules.} \\
\\ 
\texttt{Welcome to Wordle! You have six attempts to guess the target word, a valid English word of five lowercase letters (a{-}z). Please use the tags "explanation:" and "guess:" to provide a concise explanation for each guess.} \\
\\ 
\texttt{To help you make an informed guess, you will receive a clue for the word, such as} \\
\texttt{clue: snowy white.} \\
\\ 
\texttt{Here is an example guess based on the clue:} \\
\texttt{explanation: In the fairy tail Snow White, the girl is killed because she eats a poisoned apple. And the word apple has 5 letters.} \\
\texttt{guess: apple} \\
\\ 
\texttt{After each guess, your answer will be validated, and you will receive feedback indicating which letters are correct (green), which letters are correct but in the wrong position (yellow), and which letters are incorrect (red). This feedback can be useful in determining which letters to include or exclude in your next guess.} \\
\\ 
\texttt{For example, the feedback for "apple" might be:} \\
\texttt{guess\_feedback: a<yellow> p<yellow> p<green> l<yellow> e<red>} \\
\\ 
\texttt{The explanation should contain details about how both the guess\_feedback and clue work together in tandem to formulate a new guess.} \\
\\ 
\texttt{Do not generate any other text.} \\
\\ 
\texttt{Let's begin with your first guess.} \\
\\ 
\texttt{clue: winter wear} \\
            }
        }
    }
     \\ \\

    \theutterance \stepcounter{utterance}  
    & \multicolumn{4}{p{0.6\linewidth}}{
        \cellcolor[rgb]{0.561,0.737,0.561}{
            \makecell[{{p{\linewidth}}}]{
                \texttt{\tiny{[P1$\rangle$GM]}}
                \texttt{explanation: The clue "winter wear" suggests a common item like a coat, and "coats" is a valid five{-}letter word fitting the theme.} \\
\texttt{guess: coats} \\
            }
        }
    }
    & & \\ \\

    \theutterance \stepcounter{utterance}  
    & & & \multicolumn{2}{p{0.3\linewidth}}{
        \cellcolor[rgb]{0.9,0.9,0.9}{
            \makecell[{{p{\linewidth}}}]{
                \texttt{\tiny{[GM$|$GM]}}
                \texttt{attempts = 1} \\
\texttt{target\_word = glove} \\
\texttt{target\_word\_clue = winter wear} \\
\texttt{guess = coats} \\
\texttt{guess\_feedback = c<red> o<yellow> a<red> t<red> s<red>} \\
            }
        }
    }
    & & \\ \\

    \theutterance \stepcounter{utterance}  
    & & & \multicolumn{4}{p{0.6\linewidth}}{
        \cellcolor[rgb]{0.9,0.9,0.9}{
            \makecell[{{p{\linewidth}}}]{
                \texttt{\tiny{[P1$\langle$GM]}}
                \texttt{guess\_feedback: c<red> o<yellow> a<red> t<red> s<red>} \\
\\ 
\texttt{Provide your response only in this format.} \\
\texttt{explanation: details} \\
\texttt{guess: word} \\
            }
        }
    }
     \\ \\

    \theutterance \stepcounter{utterance}  
    & \multicolumn{4}{p{0.6\linewidth}}{
        \cellcolor[rgb]{0.561,0.737,0.561}{
            \makecell[{{p{\linewidth}}}]{
                \texttt{\tiny{[P1$\rangle$GM]}}
                \texttt{explanation: The feedback indicates only 'o' is correct but misplaced, while 'c', 'a', 't', 's' are incorrect; since 'o' is in the second position in "coats" but marked yellow, it must be elsewhere, and with no other letters confirmed, I will try "wools" to test 'w' and 'l' while keeping 'o' in a different spot.} \\
\texttt{guess: wool} \\
            }
        }
    }
    & & \\ \\

    \theutterance \stepcounter{utterance}  
    & & & \multicolumn{2}{p{0.3\linewidth}}{
        \cellcolor[rgb]{0.9,0.9,0.9}{
            \makecell[{{p{\linewidth}}}]{
                \texttt{\tiny{[GM$|$GM]}}
                \texttt{Error: The length of the guessed word is not 5.} \\
            }
        }
    }
    & & \\ \\

    \theutterance \stepcounter{utterance}  
    & & & \multicolumn{2}{p{0.3\linewidth}}{
        \cellcolor[rgb]{0.9,0.9,0.9}{
            \makecell[{{p{\linewidth}}}]{
                \texttt{\tiny{[GM$|$GM]}}
                \texttt{game\_result = ABORT} \\
            }
        }
    }
    & & \\ \\

    \theutterance \stepcounter{utterance}  
    & & \multicolumn{4}{p{0.6\linewidth}}{
        \cellcolor[rgb]{0.95,0.95,0.95}{
            \makecell[{{p{\linewidth}}}]{
                \texttt{\tiny{[P1$\leftrightarrow$P]}}
                \texttt{Post{-}episode Reflexion (not sent to GM):} \\
\texttt{New plan: I failed because my previous guess "wool" was only four letters long, violating the strict rule for a five{-}letter word. Additionally, my explanation was verbose and did not explicitly reference the clue "winter wear" as required. Next time, I will immediately generate a valid five{-}letter word that incorporates the confirmed 'o' (from the yellow feedback) and aligns with the "winter wear" theme, such as "scarf" or "snood," ensuring the explanation is a single concise sentence that explicitly links the clue and feedback before the guess.} \\
            }
        }
    }
    & \\ \\

\end{supertabular}
}

\end{document}
